\section{Classical theorems as two readings of one bridge}\label{sec:equivalences}

The same machinery explains why some pairs of theorems from a third-year curriculum, taught in different courses, are \q{the same theorem}: they are two readings of one invariant on one topos. Nothing in this section is new; its point is pedagogical, and it shows what a bridge does when both ends are already known.

\subsection{Equivalence relations and quotients}

Let $E$ be an equivalence relation on a set $S$, seen as a groupoid (one arrow $s\to t$ when $s\mathrel E t$). A presheaf on $E$ is a family of sets $(F_s)_{s\in S}$ with compatible bijections $F_s\cong F_t$ for $s\mathrel E t$.

\begin{itemize}
\item \emph{Deck:} the groupoid $E$ is equivalent to the discrete category $S/E$ (each connected component is contractible), so by Bridge~\ref{br:morita} $\PSh(E)\simeq\Set^{S/E}$.
\item \emph{Reading 1 (L1--L2 algebra):} families over $S$ compatible with $E$ are the same as families over $S/E$; for $F_s = \{\ast\}$ or $\emptyset$, this is the universal property of the quotient: \emph{a map $S\to\{0,1\}$ constant on classes factors through $S/E$}.
\item \emph{Reading 2 (Bridge~\ref{br:bool}):} $E$ is a groupoid, so every sub-presheaf has a complement; for subsets, \emph{a union of classes has a complement which is a union of classes}: the classes partition $S$.
\end{itemize}

\subsection{Orbits and complements}

For a group $G$, Bridge~\ref{br:bool} says that the complement of a sub-$G$-set is a sub-$G$-set. Together with the decomposition of every presheaf into connected components (true in every presheaf topos), this is the \emph{orbit decomposition}: every $G$-set is the disjoint union of its orbits, and every sub-$G$-set is a union of orbits. For a monoid that is not a group the decomposition into connected components still holds, but a component need not be \q{indecomposable} in the Boolean sense: the bridge pinpoints that the orbit decomposition theorem is a statement about invertibility.

\subsection{Birkhoff's theorem and finite $T_0$ spaces}

Let $P$ be a finite poset and $X_P$ its Alexandrov space (open sets $=$ up-sets). The deck is $\Sh(X_P)\simeq[P,\Set]$: a sheaf on $X_P$ is determined by its values on the smallest open neighbourhoods $\uparrow p$, and $\uparrow q\subseteq\uparrow p$ when $p\le q$. Two invariants of this topos:
\begin{itemize}
\item its \emph{subterminal objects} (sub-presheaves of the one-point presheaf): read on $X_P$ they are the open sets; read on $[P,\Set]$ they are the up-sets of $P$;
\item its \emph{points}: read on $X_P$ (a sober space) they are the points of $X_P$; read on $[P,\Set]$ they are the evaluations at the elements of $P$, and the principal up-sets $\uparrow p$ are exactly the join-irreducible elements of $\mathrm{Up}(P)$, with $\uparrow p\subseteq\uparrow q$ iff $q\le p$.
\end{itemize}
Reading these invariants backwards gives two classical theorems at once: \emph{Birkhoff's representation theorem} (a finite distributive lattice $L$ is isomorphic to the lattice of up-sets of the opposite of its poset of join-irreducible elements) and the \emph{correspondence between finite $T_0$ spaces and finite posets} (a finite $T_0$ space is the Alexandrov space of its specialisation order). Applying Bridge~\ref{br:bool} on top: the topos is Boolean iff $P$ is discrete, iff $\mathrm{Up}(P)$ is the power set of $P$, which is the finite case of \emph{Stone's representation theorem}: a finite Boolean algebra is the power set of its atoms.

\paragraph{A precedent in Grothendieck's notes.} In an undated manuscript note on topoi \cite{grothendieck-161-3}, p.~26, Grothendieck lists equivalent conditions on a topos $E$: to be equivalent both to the topos of a space and to a presheaf topos; to be the topos of an \q{essential} space (a space all of whose points are essential); to be equivalent to $\widehat I$ for an ordered set $I$; to be generated by the subobjects of the final object and to have enough essential points. He asks to prove that $\widehat C$ satisfies these conditions iff $C$ is preordered, and that the topos of a space $X$ does iff the essential points of $X$ are very dense, and he notes 2-equivalences between essential spatial topoi, preorders and essential spaces. These are bridges in the sense of this article, read on the deck of preorders with the arches of subterminal objects and of essential points. The generator re-derives the first exercise from the Morita arch: if $\widehat C\simeq\widehat P$ for a preorder $P$, then $\Kar(C)\simeq\Kar(P)\simeq P$, and $C$, a full subcategory of its Karoubi envelope, is thin; conversely a thin category is equivalent to its thin skeleton, a partial order (Lean: \lean{exists_preorder_moritaEquivalentCat_iff_isThin}). The second exercise, on spaces, is not formalised here.

\subsection{Cayley's theorem and the reconstruction of a group}

For a group $G$, Bridge~\ref{br:pts} says that the forgetful functor $U$ from $G$-sets to sets is, up to isomorphism, the only essential point, and that its automorphism group is $G$ (Lean: \lean{essentialPointEndMulEquivOfGroup}). Evaluating automorphisms of $U$ at the regular $G$-set $G$ gives an injective homomorphism $G\to\mathrm{Sym}(G)$: this is \emph{Cayley's theorem}. Conversely, Cayley's theorem plus the Yoneda lemma (a natural transformation of $U = \Hom(G,-)$ is determined by its value on the identity) gives the reconstruction. The two are one statement about the topos $\PSh(\B G)$, read on its points and on its representable object.

\subsection{Galois correspondence and covering spaces}

Two theorems of a third-year course, the \emph{fundamental theorem of Galois theory} (intermediate fields of a finite Galois extension $L/K$ $\leftrightarrow$ subgroups of $\mathrm{Gal}(L/K)$) and the \emph{classification of covering spaces} (connected coverings of a nice connected space $X$ $\leftrightarrow$ conjugacy classes of subgroups of $\pi_1(X)$), are two readings of one bridge, made explicit by Grothendieck \cite{sga1}: the category of finite étale $K$-algebras split by $L$, and the category of coverings of $X$, are both equivalent to a category of $G$-sets ($G = \mathrm{Gal}(L/K)$, resp.\ $\pi_1(X)$); the invariant is the poset of \emph{connected objects with a point}, which on the $G$-set side is the poset of subgroups. (These decks are not presheaf topoi on a finite shape, and this bridge is not in the Lean library; see \cite{szamuely} for an undergraduate account.)

\subsection{Confluence and extremal disconnectedness}

The last pair is in the calibration log: a preorder is confluent iff its Alexandrov space is extremally disconnected \cite[Prop.~7.10]{caramello-stone} (Lean: \lean{extremallyDisconnected_iff_confluent}). Read on a rewriting system, the Church--Rosser theorem for $\beta$-reduction becomes the extremal disconnectedness of a space of $\lambda$-terms. Besides re-deriving two known theorems, the bridge here shows that a property from programming-language theory and a property from general topology are the same property.

\begin{remark}
In each pair, one reading is usually proved in one course and the other in another, by unrelated arguments. The bridge gives no shorter proofs, but it explains why the two statements have the same shape, and it gives a recipe for finding more such pairs: pick a topos with two familiar presentations, and read an invariant on both.
\end{remark}
